% Part: computability
% Chapter: computability-theory
% Section: applications-fixed-point

\documentclass[../../../include/open-logic-section]{subfiles}

\begin{document}

\olfileid{cmp}{thy}{apf}
\olsection{স্থির-বিন্দু উপপাদ্যের প্রয়োগ}

স্থির-বিন্দু উপপাদ্য মূলত আংশিক গণনসাধ্য অপেক্ষককে তাদের সূচকের সাহায্যে
সংজ্ঞায়িত করতে দেয়। যেমন, এমন একটি সূচক~$e$ পাওয়া যায়, যাতে প্রত্যেক
$y$-এর জন্য
\[
\cfind{e}(y) = e + y.
\]
আরেকটি উদাহরণ হিসেবে, স্থির-বিন্দু উপপাদ্যের প্রমাণ ব্যবহার করে Java বা
C++-এ নিজেকেই মুদ্রণ করে এমন একটি প্রোগ্রাম তৈরি করা যায়।

মনে করো, প্রত্যেক~$e$-এর জন্য $W_e$-কে~$\cfind{e}$-এর সংজ্ঞাক্ষেত্র
ধরলে $W_0$, $W_1$, $W_2$,~\dots অনুক্রমটি গণনসাধ্যভাবে তালিকায়নযোগ্য
সেটগুলিকে তালিকায়িত করে। এই সেটগুলির কয়েকটি গণনসাধ্য। এমন কোনো
অ্যালগরিদম আছে কি, যা নিবেশ হিসেবে একটি মান~$x$ নেয় এবং $W_x$
গণনসাধ্য হলে তার চরিত্রসূচক অপেক্ষকের একটি সূচক ফেরত দেয়? উত্তর “না”;
এমন কোনো অ্যালগরিদম নেই:

\begin{thm}
নিচের বৈশিষ্ট্যসহ কোনো আংশিক গণনসাধ্য অপেক্ষক~$f$ নেই: যখনই $W_e$
গণনসাধ্য, তখন $f(e)$ সংজ্ঞায়িত এবং $\cfind{f(e)}$ তার চরিত্রসূচক অপেক্ষক।
\end{thm}

\begin{proof}
$f$-কে যেকোনো আংশিক গণনসাধ্য অপেক্ষক ধরি; আমরা এমন একটি~$e$ নির্মাণ করব,
যাতে $W_e$ গণনসাধ্য, কিন্তু $\cfind{f(e)}$ তার চরিত্রসূচক অপেক্ষক নয়।
% BN-SRC-159 TARGET-CORRECTION-BEGIN
\par\smallskip\noindent\textnormal{\small\textbf{উৎস-সংশোধন (BN-SRC-159):} উপপাদ্যটি আংশিক গণনসাধ্য প্রার্থী~$f$-কে নিষিদ্ধ করে, কিন্তু হিমায়িত ইংরেজি প্রমাণটি শুধু “যেকোনো গণনসাধ্য অপেক্ষক” ধরে, যা প্রচলিত সংজ্ঞায় সর্বত্র-সংজ্ঞায়িত। বাংলা প্রমাণে উপপাদ্যের পূর্ণ আংশিক শ্রেণিই ধরা হয়েছে।} % BN-SRC-159 TARGET-CORRECTION-END
স্থির-বিন্দু উপপাদ্য ব্যবহার করে এমন একটি সূচক~$e$ পাওয়া যায়, যাতে
\[
\cfind{e}(y) \simeq
\begin{cases}
0 & \text{যদি $y=0$ এবং $\cfind{f(e)}(0) \fdefined = 0$} \\
\text{অসংজ্ঞায়িত} & \text{অন্যথায়।}
\end{cases}
\]
অর্থাৎ, নিচে সংজ্ঞায়িত অপেক্ষকটির উপর স্থির-বিন্দু উপপাদ্য প্রয়োগ করে
$e$ পাওয়া যায়:
\[
g(x,y) \simeq
\begin{cases}
0 & \text{যদি $y=0$ এবং $\cfind{f(x)}(0) \fdefined = 0$} \\
\text{অসংজ্ঞায়িত} & \text{অন্যথায়।}
\end{cases}
\]
অনানুষ্ঠানিকভাবে $g$ যে আংশিক গণনসাধ্য তা এভাবে দেখা যায়: নিবেশ $x$ ও
$y$-এ অ্যালগরিদমটি প্রথমে পরীক্ষা করে~$y$, $0$-এর সমান কি না। সমান হলে
অ্যালগরিদমটি $f(x)$ গণনা করে, তারপর সার্বজনীন যন্ত্র ব্যবহার করে
$\cfind{f(x)}(0)$ গণনা করে। এই শেষ গণনাটি থেমে~$0$ ফেরত দিলে
অ্যালগরিদমটি~$0$ ফেরত দেয়; অন্যথায় অ্যালগরিদমটি থামে না।

প্রথমেই, প্রার্থী অপেক্ষকটি এই স্থির সূচকে অসংজ্ঞায়িত হলে নির্মিত সেটটি
শূন্য এবং তাই গণনসাধ্য, অথচ দাবিকৃত সূচকমূল্যটি সংজ্ঞায়িত নয়—সুতরাং
বৈশিষ্ট্যটি ব্যর্থ। এবার সূচকমূল্যটি সংজ্ঞায়িত ধরি। তখন
$\cfind{f(e)}(0)$ সংজ্ঞায়িত ও $0$-এর সমান হলে $\cfind{e}(y)$ ঠিক তখনই
সংজ্ঞায়িত, যখন $y$, $0$-এর সমান; তাই $W_e = \{ 0 \}$।
$\cfind{f(e)}(0)$ অসংজ্ঞায়িত হলে, অথবা সংজ্ঞায়িত কিন্তু $0$-এর সমান না
হলে, $W_e = \emptyset$। উভয় ক্ষেত্রেই $\cfind{f(e)}$, $W_e$-এর
চরিত্রসূচক অপেক্ষক নয়, কারণ নিবেশ~$0$-এ সেটি ভুল উত্তর দেয়।
% BN-SRC-160 TARGET-CORRECTION-BEGIN
\par\smallskip\noindent\textnormal{\small\textbf{উৎস-সংযোজন (BN-SRC-160):} হিমায়িত ইংরেজি যুক্তিটি $f(e)$-কে আগেই সংজ্ঞায়িত ধরে শুধু $\cfind{f(e)}(0)$-এর আচরণের দুটি ক্ষেত্র আলোচনা করে। কিন্তু নিষিদ্ধ প্রার্থী~$f$ আংশিক। বাংলা পাঠে $f(e)$ অসংজ্ঞায়িত ক্ষেত্রটি আগে নিষ্পন্ন করা হয়েছে: তখন $W_e=\emptyset$ গণনসাধ্য, অথচ প্রয়োজনীয় সূচকমূল্যই নেই।} % BN-SRC-160 TARGET-CORRECTION-END
\end{proof}

\end{document}
